\documentclass[10pt,a4paper]{amsart}
\usepackage[margin=19mm]{geometry}
\usepackage{amsmath,amssymb,mathtools}
\usepackage[scr=rsfs,cal=euler]{mathalfa}
\usepackage{xfrac}
\usepackage[unicode,hidelinks,bookmarks=false]{hyperref}
\newcommand{\ie}{i.e.\ }
\newcommand{\cf}{cf.\ }
\newcommand{\resp}{respectively\ }
\newcommand{\Subsection}{\subsection}
\providecommand{\nobreakdash}{\nobreak\hbox{-}}
\setlength{\emergencystretch}{2em}
\multlinegap0pt
\usepackage{amssymb}
\usepackage{mathtools}
\usepackage{xcolor}
\usepackage{tikz}
\usepackage{algorithm}\usepackage{algpseudocode}
\newtheorem{lemma}{Lemma}
\newcommand{\tostep}{\to}
\title{Quasilinear Equivalence Checking\\ for Detector Error Models}
\author{Mathys Rennela}
\date{Source: arXiv:2606.14677v1 (CC BY 4.0)}
\begin{document}
\maketitle

\noindent\emph{Source and licence.} Quasilinear Equivalence Checking\\ for Detector Error Models. Version arXiv:2606.14677v1 (CC BY 4.0), distributed by arXiv under the Creative Commons Attribution 4.0 International licence, http://creativecommons.org/licenses/by/4.0/, as recorded in the arXiv licence metadata for this exact version and shown on the abstract page https://arxiv.org/abs/2606.14677v1. Full licence notice: \texttt{licenses/arxiv-2606.14677-CC-BY-4.0.txt}.\par\smallskip

\noindent\textit{Editorial scope.} Counted unit: the article's lemma lem:termination (source0.tex lines 429--432). The article's complete original proof of it is delivered with it (source0.tex lines 434--451, one \texttt{proof} environment of 18 lines). No mathematics is written, repaired or re-derived: the statement and the proof are the article's own lines, in the article's own notation, and no retained line is substituted or re-flowed.  No further source-local context is needed: every cross-reference of every retained line resolves inside the two delivered spans.  Nothing was omitted from any retained span, so the package carries no omission record.  No retained line contains a citation, so the wrapper carries no bibliography and no bibliography entry.  No retained line contains a figure environment, an image-inclusion command or drawing code, so no image is delivered and the package declares no asset dependency.  Editorial formatting only, in addition: an AMS \texttt{amsart} preamble that restates the article's own declarations and macros used by the retained lines, namely the environments \texttt{lemma} on separate counters, together with the standard packages those lines need.  The retained environments print in this document as Lemma 1 in order of appearance, whereas the article prints the same items with the numbering of its own full document.  The complete native source of the article as distributed in the arXiv e-print for this exact version is bundled unchanged as \texttt{sources/202802/source0.tex}.
\begin{lemma}[Termination]
\label{lem:termination}
There is no infinite sequence of terms $t_1, t_2, \ldots$ such that $t_1 \tostep t_2 \tostep t_3 \tostep \cdots$.
\end{lemma}

\begin{proof}
Assign to each term $t$ a weight $w(t) = (C(t), T(t))$ ordered lexicographically, where $T(t)$ is the total number of targets across all error instructions in $t$ (at all scopes), and $C(t)$ is the number of nodes in the abstract syntax tree of $t$ defined by induction as follows:
\begin{align*}
  C(\varepsilon) &= 1, &
  C(\mathtt{error}(p)\,S) &= 2, \\
  C(t_1\,;\,t_2) &= C(t_1)+C(t_2)+1, &
  C(\mathtt{REPEAT}\,N\,\{t'\}) &= 1+C(t'),
\end{align*}
Rule R4 is a quotient and generates no directed reductions. Every other rule strictly decreases $w$ in lexicographic order:
\begin{itemize}
  \item R1 collapses two instruction nodes and one composition node ($C = 5$) into one instruction node ($C = 2$).
  \item R2a--R2b each replace an instruction node ($C = 2$) with $\varepsilon$ ($C = 1$).
  \item R2c--R2d remove a composition node and an $\varepsilon$ node, decreasing $C$ by $2$.
  \item R3a--R3b leave $C$ unchanged and decrease $T$ by $2$.
  \item R5 (congruence): if any step R1--R3 is applied inside the body $t$, the reduct $t'$ has strictly smaller weight $w(t') < w(t)$. Consequently $w(\mathtt{REPEAT}\,N\,\{t'\}) < w(\mathtt{REPEAT}\,N\,\{t\})$, preserving the strict decrease.
\end{itemize}
Since $w(t) \in \mathbb{N}^2$ strictly decreases at each step, no infinite rewriting sequence exists.
\end{proof}

\end{document}
